\chapter{Trabalhos Futuros}
\label{sec:future_work}

Com base nos resultados, conclusões e ameaças à validade identificadas neste estudo, delineamos várias direções promissoras para pesquisas futuras que abordam as limitações identificadas e expandem o escopo da investigação:

\section{Aprimoramentos Metodológicos}

\begin{enumerate}
\item \textbf{Otimização e Adaptação de Parâmetros:} Desenvolver estratégias para configuração automática dos parâmetros do IVF (Tamanho da Coleta $c$ e Taxa $r$), superando o viés de configuração identificado na Seção~\ref{sec:threats}. Isso inclui:
\begin{itemize}
\item Investigação de técnicas de meta-otimização para encontrar configurações ótimas específicas para cada classe de problema;
\item Desenvolvimento de mecanismos adaptativos que ajustem dinamicamente os parâmetros durante a evolução com base em métricas de diversidade populacional e convergência;
\end{itemize}

\item \textbf{Integração com Outros Algoritmos Base:} Expandir a investigação do mecanismo IVF além do SPEA2, incorporando-o sistematicamente em outros algoritmos evolutivos multiobjetivo estabelecidos (NSGA-III, MOEA/D-variants, SMS-EMOA) para avaliar sua generalidade como módulo auxiliar.
\end{enumerate}

\section{Validação Experimental Abrangente}

\begin{enumerate}
\item \textbf{Avaliação em Problemas do Mundo Real:} Abordar a limitação de benchmarks sintéticos através da aplicação sistemática do IVF/SPEA2 em problemas de otimização do mundo real de diferentes domínios, entendendo por exemplo o impacto nos exemplos reais onde o SPEA2 vem sendo aplicado.



\item \textbf{Estudos de Escalabilidade:} Investigar o comportamento do algoritmo com diferentes tamanhos de população e orçamentos computacionais, superando a limitação de configuração experimental fixa identificada nas ameaças à validade.
\end{enumerate}

\section{Extensões Algorítmicas}

\begin{enumerate}
\item \textbf{Variantes Híbridas Avançadas:} Explorar combinações do mecanismo IVF com outras técnicas de aprimoramento, como:
\begin{itemize}
\item Busca local multiobjetivo;
\item Técnicas de niching e clustering;
\item Abordagens baseadas em aprendizado de máquina para guiar a seleção.
\end{itemize}

\item \textbf{Paralelização e Computação Distribuída:} Desenvolver versões paralelas e distribuídas do IVF/SPEA2 para lidar com problemas computacionalmente intensivos e aproveitar arquiteturas de computação modernas.

\item \textbf{Integração com Técnicas de Surrogate Modeling:} Investigar a combinação do IVF/SPEA2 com modelos substitutos para problemas com avaliações de função computacionalmente caras.
\end{enumerate}

Estas direções de pesquisa futura abordam sistematicamente as limitações identificadas nas ameaças à validade e proporcionam uma base sólida para o avanço contínuo da área de otimização evolutiva multiobjetivo, contribuindo tanto para o desenvolvimento teórico quanto para aplicações práticas relevantes.